\documentclass[10pt]{article}
\usepackage[utf8]{inputenc}
\usepackage[T1]{fontenc}
\usepackage{amsmath}
\usepackage{amsfonts}
\usepackage{amssymb}
\usepackage[version=4]{mhchem}
\usepackage{stmaryrd}
\usepackage{bbold}

\begin{document}
$L 6$

\section*{L6: Congruences}
\begin{itemize}
  \item Note that Proposition 2.4 from last lecture can be stated in terms of equivalence classes under congruence modulo $m$ :\\
(a) $[a]+[c]=[a+c]$\\
(b) $[a][c]=[a c]$.
\end{itemize}

This allows us to construct addition and unultiplication tables for equivalence classes under congruence modulom using any Complete residue system modulo $m$.

$$
\begin{array}{ll|lll}
\text { Example: } & + & {[0]} & {[1]} & {[2]} \\
\hline \text { Modulo 4: } & {[0]} & {[0]} & {[1]} & {[2]} \\
& {[1]} & {[1]} & {[2]} & {[3]} \\
& {[2]} & {[2]} & {[3]} & {[0]} \\
& {[3]} & {[3]} & {[0]} & {[1]} \\
& \vdots & {[0]} & {[1]} & {[2]} \\
& {[0]} & {[0]} & {[0]} & {[0]} \\
& {[1]} & {[0]} & {[1]} & {[2]} \\
& {[2]} & {[0]} & {[2]} & {[0]} \\
& {[3]} & {[0]} & {[3]} & {[2]} \\
& [2]]
\end{array}
$$

\begin{itemize}
  \item $\mathbb{Z}_{4}:=\{[0),[1],(2),[3]\}$ along with the operations of addition and multiplication form a ring.
  \item Consider the multiplication and addition tables for equivalence classes under congruence modulo 5:
\end{itemize}

$$
\begin{array}{l|lllll}
+ & {[0]} & {[1]} & {[2]} & {[3]} & {[4]} \\
\hline[0] & {[0]} & {[1]} & {[2]} & {[3]} & {[4]} \\
{[1]} & {[1]} & {[2]} & {[3]} & {[4]} & {[0]} \\
{[2]} & {[2]} & {[3]} & {[4]} & {[0]} & {[1]} \\
{[3]} & {[3]} & {[4]} & {[0]} & {[1]} & {[2]} \\
{[4]} & {[4]} & {[0]} & {[1]} & {[2]} & {[3]} \\
{[0]} & {[0]} & {[1]} & {[2]} & {[3]} & {[4]} \\
\hline[0] & {[0]} & {[0]} & {[0]} & {[0]} \\
{[1]} & {[0]} & {[1]} & {[2]} & {[3]} & {[4]} \\
{[2]} & {[0]} & {[2]} & {[4]} & {[1]} & {[3]} \\
{[3]} & {[0]} & {[3]} & {[1]} & {[4]} & {[2]} \\
{[4]} & {[0]} & {[4]} & {[3]} & {[2]} & {[1]}
\end{array}
$$

\begin{itemize}
  \item $\mathbb{Z}_{5}:=\{[0],[1],[2],[3],[4]\}$ along with the operations of addition and multiplication form a field.
\end{itemize}

Question: Why does $\mathbb{Z}_{5}$ have more structure than $\mathbb{Z}_{4}$ ?

\begin{itemize}
  \item Cancellation law for multiplication does not hold in general for congruence modulo $m$.\\
Proposition 2.5:\\
Let $a, b, c \in \mathbb{Z}$. Then $c a \equiv c b(\bmod m)$ if and only if $a \equiv b \bmod (m /(c, m))$.\\
Proof:\\
$\Leftrightarrow$ direction): Assume that $c a \equiv c b(\bmod m)$ and let $d:=(c, m)$. We have $m \mid c a-c b=c(a-b)$ and thus $\mathrm{m} / \mathrm{d} \mid(c / d)(a-b)$, and so $a \equiv b(\bmod \mathrm{~m} / \mathrm{d})$ Since $(\mathrm{m} / \mathrm{d}, c / \mathrm{d})=1$ by Proposition 1.10.\\
$(\Leftarrow$ direction $):$ Let $d:=(c, m)$ and assume that $a \equiv b(\bmod m / d)$. Then $m / d \mid a-b$, so that implies that $m \mid d a-d b$. So $m \mid(4 d)(d a-d b)=c a-c b$\\
and $c a \equiv c b(\bmod m)$, as desired.
\end{itemize}

\section*{Linear congruences in one variable:}
Definition: Let $a, b \in \mathbb{Z}$. A congruence of the\\
form $a x \equiv b(\bmod m)$ is said to be a linear congrnence in the variable $x$.\\
Example: $2 x \equiv 3(\bmod 4)$ has no solutions in $\mathbb{Z}$. Check $x \equiv 0,1,2,3(\bmod 4)$ don't work.

\begin{itemize}
  \item $2 x \equiv 3(\bmod 5)$ has one solution modulo 5 , namely $x \equiv 4(\bmod 5)$, and so has solution Set in $\mathbb{Z}$ given by $[4]$.
  \item $2 x \equiv 4(\bmod b)$ has two solutions modulo b, namely $x \equiv 2,5(\bmod 6)$, and so has solution Set in $\mathbb{Z}$ given by [2] U[5].
  \item $3 x \equiv 9(\bmod 6)$ has three solutions modulo 6) namely $x \equiv 1,3,5(\bmod 6)$, and so has solution set in $\mathbb{Z}$ given by $[1] \cup[3] \cup[5]$.
\end{itemize}

Theorem 2.6: Let $a x \equiv b(\operatorname{modm})$,\\
and let $d=(a, m)$. If $d+b$, then the congruence has no solutions in $\mathbb{Z}$. If $d \mid b$, then the congruence has exactly $d$ incongruent solutions modulo $m$ in $\mathbb{Z}$.\\
Proof: Observe that $a x \equiv b \bmod m)$ if and only if $m / a x-b$ if and only if $a x-b=m y$ for some $y \in \mathbb{Z}$ if and only if $a x-m y=b$ for some $y \in \mathbb{Z}$. Thus $a x \equiv b(\bmod m)$ is Solvable in $x$ if and only if $a x-m y=b$ is solvable for $x$ and $y(*)$.\\
Since $d \mid a$ and $d \mid m$, we have that $d \mid a x-m y$ by Proposition 1.2, which implies that $d \| b$. Equivalently, if $d+b$, then $a x \equiv b(\bmod m)$ has no solutions. Wlog now assume that $d l b$.\\
We now prove the Second part of the Thervem in 4 steps:\\
(a) We show that $a x \equiv b(\bmod m)$ has $a$ solution, say $x_{0}$.\\
(b) Given any solution $x_{0}$, we show that\\
$a x \equiv b(\bmod m)$ has infinitely many solutions in $\mathbb{Z}$ of a given form.\\
(c) Given any solution $x_{0}$, we show that every solution to $a x \equiv b(\bmod m)$ has the form in (b) (from which (b) and (c) combine to yield all solutions in $\mathbb{Z}$ to $a x \equiv b(\bmod m))$.\\
(d) We show that there are exactly $d$ incongruent solutions modulo $m$ in $\mathbb{Z}$.\\
For (a): By Proposition 1.11, there exist $r, s \in \mathbb{Z}$ such that $d=(a, m)=a r+m s$.\\
Furthermore, $d 1 b$ implies that $b=d e$ for some $e \in \mathbb{Z}$. Then

$$
b=d e=(a r+m s) e=a(r e)+m(s e) .
$$

Thus $x=$ re and $y=-$ se solve $a x-m y=b$,\\
from which $x_{0}=$ re solves $a x \equiv b$ (modm)\\
by (*).

For (b): Let $x_{0}$ be any solution of $a x \equiv b(\bmod m)$. Let $n \in \mathbb{Z}$ and consider $x_{0}+\left(\frac{m}{d}\right) n$. Note that $x_{0}+\left(\frac{m}{d}\right) n$ is an integer, since $d / m$. Furthermore,

$$
\begin{aligned}
a\left(x_{0}+\left(\frac{m}{d}\right) n\right) & =a x_{0}+a\left(\frac{m}{d}\right) n \\
& =a x_{0}+\left(\frac{a}{d}\right) m n \\
& =b+0(\bmod m) \\
& \equiv b(\bmod m) .
\end{aligned}
$$

Thus given any solution $x_{0}$ to $a x \equiv b(\bmod m)$, the infinitely many values $x_{0}+\left(\frac{y}{d}\right) n$ in $\mathbb{Z}$, $n \in \mathbb{Z}$, are also solutions.\\
For (c) : Let $x_{0}$ be any solution to $a x \equiv b(\bmod m)$. Then $a x_{0}-m y_{0}=b$ for some $y_{0} \in \mathbb{Z}$ by\\
Any other solution $x$ of $a x \equiv b$ (modm) implies the existence of $y \in \mathbb{Z}$ such that $a x-m y=b$ by (*), and so we have

$$
(a x-m y)-\left(a x_{0}-m y_{0}\right)=b-b=0 .
$$

Thus $a\left(x-x_{0}\right)=m\left(y-y_{0}\right)$.

Dividing both sides by $d:=(a, m)$ we obtain $\frac{a}{d}\left(x-x_{0}\right)=\frac{m}{d}\left(y-y_{0}\right)$.\\
Now $\frac{m}{d} \left\lvert\,\left(\frac{a}{d}\right)\left(x-x_{0}\right)\right.$; since $\left(\frac{a}{d}, \frac{m}{d}\right)=1$ by Proposition 1.10, we have that $\left.\frac{m}{d} \right\rvert\, x-x_{0}$. Consequently $x-x_{0}=\left(\frac{m}{d}\right)_{n}$ for some $n \in \mathbb{Z}$,\\
So equivalently, $x=x_{0}+\left(\frac{m}{d}\right) n$ for some $n \in \mathbb{Z}$. Hence, given a solution $x_{0}$ of $a x \equiv b(\operatorname{modm})$, we have that every solution takes the form $x_{0}+\left(\frac{m}{d}\right) n$ with $n \in \mathbb{Z}$. Combining this vesult with part (b), we have that all solutions are given precisely by $x_{0}+\left(\frac{y}{d}\right)_{n}$ with $n \in \mathbb{Z}$.\\
For (d): We find necessary and sufficient conditions for the congruence modulo m of two solutions in the family $x_{0}+\left(\frac{m}{d}\right) n, n \in \mathbb{Z}$.

We have

$$
x_{0}+\left(\frac{m}{d}\right) n_{1} \equiv x_{0}+\left(\frac{m}{d}\right) n_{2}(\bmod m)
$$

if and only if

$$
\left(\frac{m}{d}\right) n_{1} \equiv\left(\frac{m}{d}\right) n_{2}(\bmod m)
$$

if and only if

$$
m \left\lvert\,\left(\frac{m}{d}\right)\left(n_{1}-n_{2}\right)\right.
$$

if and only if (by Proposition 2.5)

$$
d \mid n_{1}-n_{2}
$$

if and only if

$$
n_{1} \equiv n_{2}(\bmod d) .
$$

Thus two solutions of the form $x_{0}+\left(\frac{m}{d}\right) n$, $n \in \mathbb{Z}$, are congruent modulo $m$ if and only if the corresponding $n$-values are congrwent modulo $d$. Thus a complete set of incongrnent Solutions modulo $m$ of $a x \equiv b(\bmod m)$ is obtained from $x_{0}+\left(\frac{m}{d}\right)^{n}, n \in \mathbb{Z}$, by letting n range over a complete residue system modulo $d$, say $\{0,1, \ldots, d-1\}$ for example.

Corollary 2.7: Consider $a x \equiv b(\bmod m)$\\
and let $d=(a, m)$. If $d \| b$, then there are $d$ incongruent solutions modulo $m$, given by

$$
x_{0}+\left(\frac{m}{d}\right) n, \quad n=0,1,2, \ldots, d-1,
$$

where $x_{0}$ is any particular solution to the Congruence.\\
Example: Find all incongruent solutions to $16 x \equiv 8(\bmod 28)$ Lets use the Euclidean algorithm to Compute

$$
\begin{aligned}
d=(16,28) ; & 28=(6.1+12 \\
& 16=12.1+4 \\
& 12=4.3+0,
\end{aligned}
$$

So $d=4$. Since 418 , the Congruence has four incongruent solutions modulo 28 by\\
Theorem 2.6. Working backward through the\\
Euclidean algorithm we obtain

$$
4=2 \cdot 16+(-1) 28 .
$$

Thus $8=4 \cdot 16+(-2) \cdot 28$,\\
and so $x_{0}=4$ is a particular solution to $16 x \equiv 8(\bmod 28)$. Thus by Corollary 2.7, all incongruent solutions to $16 x \equiv 8(\bmod 28)$ are given by

$$
4+\left(\frac{28}{4}\right) n, n=0,1,2,3,
$$

or equivalently, $x \equiv 4,11,18,25(\bmod 28)$.\\
Definition: Any solution of $a x \equiv 1(\bmod m)$ is said to be the multiplicative inverse of $a$ modulo m.

Corollary 2.8 (to Theovem 2.6): The Congruence $a x \equiv 1(\bmod m)$ has a solution if and only if $(a, m)=1$.\\
If $(a, m)=1$, then $a x \equiv 1(\bmod m)$ has exactly one incongruent solution modulo $m$.\\
Example: Find all incongruent solutions to $5 x \equiv 12(\bmod 16)$. Since $(5,16)=1$, the congmence has a unique solution modulo 16 by Theovem 2.6. Use the method of the\\
previous example (exercise) to show that 13 is the inverse of 5 modulo 16. Thus

$$
13 \cdot 5 x \equiv 13 \cdot 12(\bmod 16)
$$

and so

$$
\begin{aligned}
x & \equiv 12.13(\bmod 16) \\
& \equiv(-4)(-3)(\bmod 16) \\
& \equiv 12(\bmod 16) .
\end{aligned}
$$


\end{document}